\chapter{내적공간과 직교화}

\chapterintro{이 장에서는 벡터공간에 길이와 각도를 도입하여 기하적 해석이 가능한 선형대수의 핵심 틀을 정리합니다. 내적의 공리에서 출발하여 Cauchy--Schwarz 부등식, Gram--Schmidt 직교화, 직교투영, 최소제곱 문제까지 한 흐름으로 연결하겠습니다. 마지막에는 QR 분해를 통해 계산 절차를 안정적으로 구현하는 방법까지 확인하시게 됩니다.}
\chaptergoals{내적공간의 기본 정의와 핵심 부등식을 엄밀히 증명할 수 있다.}{Gram--Schmidt 과정을 이용해 정규직교기저를 구성하고 좌표를 계산할 수 있다.}{직교투영과 최소제곱의 구조를 정상방정식과 QR 분해로 설명할 수 있다.}

\sectiontemplate{내적, 노름, 기본 부등식}{내적(inner product)을 도입하면 벡터공간에서 길이와 각도를 동시에 다룰 수 있습니다. 이 절에서는 내적의 공리를 정확히 고정하고, 이후 모든 절에서 반복해서 쓰이는 기본 부등식을 완전증명으로 확보하겠습니다.}{\item 내적의 공리와 표기 약속을 일관되게 사용할 수 있습니다.\item Cauchy--Schwarz 부등식과 삼각부등식을 증명할 수 있습니다.\item 노름, 거리, 직교의 관계를 계산 문제에 적용할 수 있습니다.}{\item 복소수 켤레와 절댓값\item 벡터공간의 선형결합\item 실수부 $\operatorname{Re}(z)$의 성질}

\begin{definition}
$V$를 $\R$ 또는 $\C$ 위의 벡터공간이라 하자. 사상
$$
\langle\cdot,\cdot\rangle:V\times V\to\C
$$
가 다음 조건을 만족하면 $V$의 \emph{내적}(inner product)이라 한다.
\begin{enumerate}[label=(\roman*)]
\item $\langle ax_1+bx_2,y\rangle=a\langle x_1,y\rangle+b\langle x_2,y\rangle$ \quad ($a,b\in\C$ 혹은 $\R$),
\item $\langle x,y\rangle=\overline{\langle y,x\rangle}$,
\item $\langle x,x\rangle\ge 0$,
\item $x\ne0$이면 $\langle x,x\rangle>0$.
\end{enumerate}
\end{definition}

\begin{definition}
내적공간 $V$에서
$$
\|x\|:=\sqrt{\langle x,x\rangle},\qquad d(x,y):=\|x-y\|
$$
를 각각 \emph{노름}(norm), \emph{거리}(distance)라 한다. 또한 $\langle x,y\rangle=0$이면 $x,y$가 \emph{직교}(orthogonal)한다고 한다.
\end{definition}

\begin{theorem}[Cauchy--Schwarz 부등식]
임의의 $x,y\in V$에 대하여
$$
|\langle x,y\rangle|\le \|x\|\,\|y\|
$$
가 성립한다.
\end{theorem}

\paragraph{증명 전략.}
핵심은 항상 음이 아닌 양수인 양 $\|x-\lambda y\|^2$를 적절한 $\lambda$에서 평가하는 것입니다. $\lambda$를 내적값으로 선택하면 교차항이 정확히 소거되어 원하는 부등식이 곧바로 나옵니다.

\begin{proof}
$y=0$이면 $\langle x,y\rangle=0$이므로 성립한다. 이제 $y\ne0$이라 하자. $\langle y,y\rangle>0$이고 실수이므로
$$
\lambda:=\frac{\langle x,y\rangle}{\langle y,y\rangle}
$$
를 정의할 수 있다. 양의정부호성에 의해
$$
0\le \|x-\lambda y\|^2
= \langle x-\lambda y,x-\lambda y\rangle.
$$
내적의 성질을 전개하면
\begin{align*}
\langle x-\lambda y,x-\lambda y\rangle
&=\langle x,x\rangle-\langle x,\lambda y\rangle-\langle \lambda y,x\rangle+\langle \lambda y,\lambda y\rangle\\
&=\|x\|^2-\overline{\lambda}\langle x,y\rangle-\lambda\langle y,x\rangle+|\lambda|^2\langle y,y\rangle.
\end{align*}
또한 $\langle y,x\rangle=\overline{\langle x,y\rangle}$이므로
$$
\overline{\lambda}\langle x,y\rangle
=\frac{|\langle x,y\rangle|^2}{\langle y,y\rangle},\qquad
\lambda\langle y,x\rangle
=\frac{|\langle x,y\rangle|^2}{\langle y,y\rangle},
$$
$$
|\lambda|^2\langle y,y\rangle
=\frac{|\langle x,y\rangle|^2}{\langle y,y\rangle}.
$$
따라서
$$
0\le \|x\|^2-\frac{|\langle x,y\rangle|^2}{\langle y,y\rangle}
=\|x\|^2-\frac{|\langle x,y\rangle|^2}{\|y\|^2}.
$$
양변에 $\|y\|^2$를 곱하면
$$
|\langle x,y\rangle|^2\le \|x\|^2\|y\|^2.
$$
양변이 음이 아니므로 제곱근을 취해 결론을 얻는다.
\end{proof}

\paragraph{의미 해석.}
이 정리는 내적이 만드는 ``각도''가 항상 유한하고 안정적임을 보장합니다. 이후 직교투영에서 계수 $\langle x,e_i\rangle$를 사용할 수 있는 이유도 이 부등식 덕분에 계산량이 폭주하지 않기 때문입니다.

\begin{theorem}[삼각부등식]
임의의 $x,y\in V$에 대하여
$$
\|x+y\|\le \|x\|+\|y\|
$$
가 성립한다.
\end{theorem}

\paragraph{증명 전략.}
$\|x+y\|^2$를 내적으로 전개한 뒤 교차항을 Cauchy--Schwarz로 제어합니다. 마지막에 완전제곱 형태로 묶어 제곱근을 취하면 결과가 나옵니다.

\begin{proof}
내적 전개로
\begin{align*}
\|x+y\|^2
&=\langle x+y,x+y\rangle\\
&=\|x\|^2+\langle x,y\rangle+\langle y,x\rangle+\|y\|^2\\
&=\|x\|^2+2\operatorname{Re}\langle x,y\rangle+\|y\|^2.
\end{align*}
실수부의 절댓값 부등식과 Cauchy--Schwarz를 적용하면
$$
2\operatorname{Re}\langle x,y\rangle
\le 2|\langle x,y\rangle|
\le 2\|x\|\|y\|.
$$
따라서
$$
\|x+y\|^2\le \|x\|^2+2\|x\|\|y\|+\|y\|^2=(\|x\|+\|y\|)^2.
$$
양변이 음이 아니므로 제곱근을 취하면
$$
\|x+y\|\le \|x\|+\|y\|.
$$
\end{proof}

\paragraph{의미 해석.}
삼각부등식은 노름이 실제 ``거리''처럼 동작함을 보여 줍니다. 최단거리, 근사해, 오차경계 같은 해석이 모두 이 성질 위에서 전개됩니다.

\begin{example}
$\R^3$의 표준내적 $\langle x,y\rangle=x_1y_1+x_2y_2+x_3y_3$를 생각하자.
$$
x=(1,2,-1),\qquad y=(2,0,1)
$$
이면
$$
\langle x,y\rangle=1,\qquad
\|x\|=\sqrt6,\qquad
\|y\|=\sqrt5.
$$
따라서
$$
|\langle x,y\rangle|=1\le \sqrt{30}=\|x\|\,\|y\|
$$
가 확인된다.
\end{example}

\subsection*{주의}
\begin{warning}
복소수 내적에서는 어느 변수가 선형이고 어느 변수가 켤레선형인지 약속을 먼저 고정해야 한다. 약속이 바뀌면 Gram--Schmidt 계수식도 함께 바뀐다.
\end{warning}
\begin{warning}
$\|x\|^2=\langle x,x\rangle$는 항상 실수이지만, $\langle x,y\rangle$는 복소수가 될 수 있다. 실수부와 절댓값을 구분하지 않으면 부등식 전개에서 오류가 난다.
\end{warning}

\subsection*{자가진단퀴즈}
\begin{enumerate}[label=\arabic*.]
\item $\C^2$에서 $\langle x,y\rangle=x_1\overline{y_1}+x_2\overline{y_2}$가 내적임을 공리로 확인하라.
\item Cauchy--Schwarz 부등식의 등호가 언제 성립하는지 서술하고, 한 쌍의 예를 제시하라.
\item 삼각부등식 증명에서 $\operatorname{Re}\langle x,y\rangle\le |\langle x,y\rangle|$가 쓰이는 위치를 정확히 찾고 이유를 설명하라.
\end{enumerate}

\sectiontemplate{Gram--Schmidt 직교화와 정규직교기저}{직교기저를 갖추면 좌표 계산과 투영 계산이 동시에 단순해집니다. 이 절에서는 Gram--Schmidt 직교화 과정을 완전증명으로 정리하고, 정규직교기저에서 좌표를 읽는 공식을 확보하겠습니다.}{\item Gram--Schmidt 알고리즘의 각 단계가 왜 정당한지 증명할 수 있습니다.\item 유한차원 내적공간에서 정규직교기저를 구성할 수 있습니다.\item 정규직교기저 좌표식과 Parseval 항등식을 사용할 수 있습니다.}{\item 선형독립과 기저\item Cauchy--Schwarz 부등식\item 부분공간의 차원}

\begin{definition}
벡터열 $(e_1,\dots,e_n)$이
$$
\langle e_i,e_j\rangle=
\begin{cases}
1,& i=j,\\
0,& i\ne j
\end{cases}
$$
를 만족하면 \emph{정규직교집합}(orthonormal set)이라 한다. 또한 정규직교집합이 공간 $V$를 생성하면 \emph{정규직교기저}(orthonormal basis)라 한다.
\end{definition}

\begin{theorem}[Gram--Schmidt 직교화]
내적공간 $V$에서 선형독립인 벡터열 $v_1,\dots,v_n$을 잡자. 다음과 같이 정의한다.
$$
u_1:=v_1,\qquad
u_k:=v_k-\sum_{j=1}^{k-1}\frac{\langle v_k,u_j\rangle}{\langle u_j,u_j\rangle}u_j\quad(2\le k\le n).
$$
그러면 다음이 성립한다.
\begin{enumerate}[label=(\roman*)]
\item $u_1,\dots,u_n$은 서로 직교하고 각 $u_k\ne0$이다.
\item 모든 $k$에 대해 $\operatorname{span}(u_1,\dots,u_k)=\operatorname{span}(v_1,\dots,v_k)$이다.
\item $e_k:=u_k/\|u_k\|$로 두면 $e_1,\dots,e_n$은 정규직교집합이며 같은 부분공간을 생성한다.
\end{enumerate}
\end{theorem}

\paragraph{증명 전략.}
귀납법으로 ``이미 만든 벡터들은 직교하고 같은 부분공간을 생성한다''는 불변량을 유지합니다. 새 벡터에서 이전 방향 성분을 모두 제거하면 직교성이 생기고, 영벡터가 되지 않음을 선형독립성으로 확인합니다.

\begin{proof}
$k=1$에서는 자명하다. 이제 $k\ge2$에 대해 귀납 가정으로 $u_1,\dots,u_{k-1}$이 서로 직교하고
$$
\operatorname{span}(u_1,\dots,u_{k-1})
=\operatorname{span}(v_1,\dots,v_{k-1})
$$
라 하자.

먼저 직교성을 보이자. $m<k$에 대해
\begin{align*}
\langle u_k,u_m\rangle
&=\left\langle v_k-\sum_{j=1}^{k-1}\frac{\langle v_k,u_j\rangle}{\langle u_j,u_j\rangle}u_j,\ u_m\right\rangle\\
&=\langle v_k,u_m\rangle-\sum_{j=1}^{k-1}\frac{\langle v_k,u_j\rangle}{\langle u_j,u_j\rangle}\langle u_j,u_m\rangle.
\end{align*}
$u_j$들이 서로 직교하므로 $\langle u_j,u_m\rangle=0$ ($j\ne m$)이고 $j=m$일 때만 남는다. 따라서
$$
\langle u_k,u_m\rangle
=\langle v_k,u_m\rangle
-\frac{\langle v_k,u_m\rangle}{\langle u_m,u_m\rangle}\langle u_m,u_m\rangle
=0.
$$
즉 $u_k$는 이전 벡터들과 모두 직교한다.

다음으로 $u_k\ne0$임을 보이자. 만약 $u_k=0$이면
$$
v_k=\sum_{j=1}^{k-1}\frac{\langle v_k,u_j\rangle}{\langle u_j,u_j\rangle}u_j
\in \operatorname{span}(u_1,\dots,u_{k-1})
=\operatorname{span}(v_1,\dots,v_{k-1}),
$$
이는 $v_1,\dots,v_n$의 선형독립성에 모순이다.

부분공간 일치는 다음과 같다. 정의식으로 $u_k\in\operatorname{span}(v_1,\dots,v_k)$이므로
$$
\operatorname{span}(u_1,\dots,u_k)\subset \operatorname{span}(v_1,\dots,v_k).
$$
반대로 정의식을 정리하면
$$
v_k=u_k+\sum_{j=1}^{k-1}\frac{\langle v_k,u_j\rangle}{\langle u_j,u_j\rangle}u_j
\in\operatorname{span}(u_1,\dots,u_k).
$$
그리고 귀납 가정에 의해 $v_1,\dots,v_{k-1}\in\operatorname{span}(u_1,\dots,u_k)$이므로 반대 포함도 성립한다.

귀납을 마치면 (i), (ii)가 모두 성립한다. 마지막으로 $e_k=u_k/\|u_k\|$는 $\|u_k\|>0$이므로 잘 정의되고,
$$
\langle e_i,e_j\rangle
=\frac{\langle u_i,u_j\rangle}{\|u_i\|\,\|u_j\|}
=\delta_{ij}
$$
이므로 정규직교집합이며, 스칼라배만 바뀌었으므로 생성 부분공간은 동일하다.
\end{proof}

\paragraph{의미 해석.}
이 정리는 ``임의의 독립 기저를 계산 친화적인 직교기저로 바꾸는 알고리즘''을 제공합니다. 이후 투영, 최소제곱, QR 분해는 모두 이 알고리즘을 핵심 부품으로 사용합니다.

\begin{theorem}[정규직교기저 좌표식과 Parseval 항등식]
$V$가 유한차원 내적공간이고 $\{e_1,\dots,e_n\}$이 정규직교기저라 하자. 임의의 $x\in V$에 대하여
$$
x=\sum_{i=1}^n \langle x,e_i\rangle e_i
$$
가 성립하며, 또한
$$
\|x\|^2=\sum_{i=1}^n |\langle x,e_i\rangle|^2
$$
가 성립한다.
\end{theorem}

\paragraph{증명 전략.}
후보 표현 $x-\sum_i\langle x,e_i\rangle e_i$를 만들어 모든 기저벡터와 직교함을 보입니다. 기저 전체와 직교하면 영벡터라는 사실로 좌표식을 얻고, 이를 다시 내적에 대입하여 노름 공식을 계산합니다.

\begin{proof}
다음을 두자.
$$
y:=x-\sum_{i=1}^n \langle x,e_i\rangle e_i.
$$
각 $j$에 대해
\begin{align*}
\langle y,e_j\rangle
&=\langle x,e_j\rangle-\sum_{i=1}^n \langle x,e_i\rangle\langle e_i,e_j\rangle\\
&=\langle x,e_j\rangle-\langle x,e_j\rangle=0.
\end{align*}
즉 $y$는 모든 $e_j$와 직교한다. 그런데 $\{e_1,\dots,e_n\}$은 $V$의 기저이므로 $y$는 $V$의 모든 벡터와 직교한다. 특히 자기 자신과도 직교하므로
$$
\langle y,y\rangle=0.
$$
양의정부호성에 의해 $y=0$이고, 따라서
$$
x=\sum_{i=1}^n \langle x,e_i\rangle e_i.
$$

이제 $c_i:=\langle x,e_i\rangle$라 두고 위 식을 내적에 대입하면
\begin{align*}
\|x\|^2
&=\left\langle \sum_{i=1}^n c_i e_i,\sum_{j=1}^n c_j e_j\right\rangle\\
&=\sum_{i=1}^n\sum_{j=1}^n c_i\overline{c_j}\langle e_i,e_j\rangle\\
&=\sum_{i=1}^n |c_i|^2
=\sum_{i=1}^n |\langle x,e_i\rangle|^2.
\end{align*}
\end{proof}

\paragraph{의미 해석.}
정규직교기저에서는 좌표를 구하기 위해 연립방정식을 풀 필요가 없습니다. 내적 한 번으로 좌표가 바로 나오며, 에너지(노름 제곱)가 좌표 제곱합으로 분해됩니다.

\begin{example}
$\R^3$에서
$$
v_1=(1,1,0),\quad v_2=(1,-1,0),\quad v_3=(1,0,1)
$$
을 잡자. $v_1,v_2,v_3$는 선형독립이다. Gram--Schmidt를 적용하면
$$
u_1=v_1,\qquad
u_2=v_2,\qquad
u_3=v_3-\frac{\langle v_3,u_1\rangle}{\langle u_1,u_1\rangle}u_1-\frac{\langle v_3,u_2\rangle}{\langle u_2,u_2\rangle}u_2=(0,0,1).
$$
따라서 정규화된 기저는
$$
e_1=\frac1{\sqrt2}(1,1,0),\quad
e_2=\frac1{\sqrt2}(1,-1,0),\quad
e_3=(0,0,1)
$$
이 된다.
\end{example}

\subsection*{주의}
\begin{warning}
Gram--Schmidt 결과는 입력 순서에 의존한다. 같은 부분공간이라도 벡터 순서를 바꾸면 얻어지는 정규직교기저가 달라질 수 있다.
\end{warning}
\begin{warning}
중간 단계에서 $u_k=0$이 나오면 계산 실수가 아니라 원래 벡터열의 선형종속을 뜻한다. 이때는 독립한 부분열만 추려서 다시 시작해야 한다.
\end{warning}

\subsection*{자가진단퀴즈}
\begin{enumerate}[label=\arabic*.]
\item $\R^3$의 벡터열 $(1,1,1),(1,1,0),(1,0,0)$에 Gram--Schmidt를 적용하여 정규직교기저를 구하라.
\item 정규직교기저에서 좌표식 $x=\sum_i\langle x,e_i\rangle e_i$가 왜 선형대수적 좌표식과 일치하는지 설명하라.
\item Parseval 항등식이 Cauchy--Schwarz 부등식을 다시 얻는 데 어떻게 쓰이는지 서술하라.
\end{enumerate}

\sectiontemplate{직교여공간, 직교분해, 직교투영}{부분공간에 대한 직교여공간을 도입하면 ``벡터를 부분공간 성분과 오차 성분으로 나누는'' 표준 분해가 가능해집니다. 이 절에서 분해의 존재와 유일성을 증명하고, 최소거리 문제와 직교투영의 동치를 확인하겠습니다.}{\item 직교여공간의 정의와 기본 성질을 사용할 수 있습니다.\item $V=W\oplus W^\perp$ 정리를 완전증명으로 설명할 수 있습니다.\item 직교투영이 최소거리점과 동치임을 증명할 수 있습니다.}{\item 정규직교기저\item 부분공간과 직합\item Pythagoras 항등식}

\begin{definition}
내적공간 $V$의 부분공간 $W\subset V$에 대하여
$$
W^\perp:=\{v\in V:\langle v,w\rangle=0\ \text{for all }w\in W\}
$$
를 $W$의 \emph{직교여공간}(orthogonal complement)이라 한다.
\end{definition}

\begin{theorem}[직교분해 정리]
$V$가 유한차원 내적공간이고 $W\subset V$가 부분공간이면
$$
V=W\oplus W^\perp
$$
가 성립한다. 특히
$$
\dim W^\perp=\dim V-\dim W.
$$
\end{theorem}

\paragraph{증명 전략.}
먼저 $W$의 정규직교기저를 잡고 이를 $V$의 정규직교기저로 확장합니다. 임의의 벡터를 앞부분(=$W$ 성분)과 뒷부분(=$W^\perp$ 성분)으로 분리하면 존재가 나오고, 교집합이 영벡터뿐임을 보여 유일성을 얻습니다.

\begin{proof}
$\dim W=k$, $\dim V=n$이라 하자. $W$의 기저를 Gram--Schmidt로 직교화하여 $W$의 정규직교기저
$$
e_1,\dots,e_k
$$
를 얻는다. 이를 $V$의 정규직교기저
$$
e_1,\dots,e_k,e_{k+1},\dots,e_n
$$
로 확장할 수 있다.

임의의 $v\in V$에 대해
$$
w:=\sum_{i=1}^k\langle v,e_i\rangle e_i,\qquad
u:=v-w
$$
로 둔다. 정의상 $w\in W$이다. 또한 $1\le j\le k$에 대해
\begin{align*}
\langle u,e_j\rangle
&=\left\langle v-\sum_{i=1}^k\langle v,e_i\rangle e_i,\ e_j\right\rangle\\
&=\langle v,e_j\rangle-\sum_{i=1}^k\langle v,e_i\rangle\langle e_i,e_j\rangle\\
&=0.
\end{align*}
$W$의 임의의 원소 $w'\in W$는 $e_1,\dots,e_k$의 선형결합이므로 $\langle u,w'\rangle=0$이다. 따라서 $u\in W^\perp$이고
$$
v=w+u\in W+W^\perp.
$$
$v$가 임의였으므로 $V=W+W^\perp$.

이제 교집합을 보자. $z\in W\cap W^\perp$이면 특히 $z\in W$이므로 직교조건에서 $w=z$를 대입하여
$$
\langle z,z\rangle=0.
$$
양의정부호성으로 $z=0$이다. 따라서 $W\cap W^\perp=\{0\}$이고
$$
V=W\oplus W^\perp.
$$
직합의 차원 공식으로
$$
\dim V=\dim W+\dim W^\perp
$$
이므로 결론이 따른다.
\end{proof}

\paragraph{의미 해석.}
모든 벡터는 ``모형 성분($W$) + 오차 성분($W^\perp$)''으로 유일하게 분해됩니다. 최소제곱 해석에서 잔차가 왜 열공간과 직교해야 하는지가 여기서 이미 결정됩니다.

\begin{definition}
위 정리의 분해
$$
v=w+u,\qquad w\in W,\ u\in W^\perp
$$
에서 $w$를 $v$의 $W$ 위 \emph{직교투영}(orthogonal projection)이라 하며 $P_W(v)=w$로 쓴다.
\end{definition}

\begin{theorem}[최소거리 정리]
$V$가 유한차원 내적공간, $W\subset V$가 부분공간이라 하자. $v\in V$, $w_0\in W$에 대하여 다음 두 조건은 동치이다.
\begin{enumerate}[label=(\roman*)]
\item $v-w_0\in W^\perp$.
\item 모든 $w\in W$에 대해 $\|v-w_0\|\le \|v-w\|$.
\end{enumerate}
또한 이 조건을 만족하는 $w_0$는 유일하다.
\end{theorem}

\paragraph{증명 전략.}
(i)$\Rightarrow$(ii)는 직교분해에 대한 Pythagoras 항등식으로 즉시 보입니다. (ii)$\Rightarrow$(i)는 최소성 가정 아래 $w_0$를 미소하게 이동한 벡터 $w_0+tz$를 대입하여, 잔차와 $W$의 임의 방향 $z$의 내적이 0이어야 함을 강제로 얻습니다.

\begin{proof}
먼저 (i)를 가정하자. 임의의 $w\in W$에 대해
$$
v-w=(v-w_0)+(w_0-w).
$$
여기서 $v-w_0\in W^\perp$, $w_0-w\in W$이므로 두 벡터는 직교한다. 따라서
$$
\|v-w\|^2=\|v-w_0\|^2+\|w-w_0\|^2\ge \|v-w_0\|^2,
$$
즉 $\|v-w_0\|\le\|v-w\|$이다.

다음으로 (ii)를 가정하자. $r:=v-w_0$라 두고 임의의 $z\in W$를 택한다. 모든 $t\in\R$ 또는 $\C$에 대해 $w_0+tz\in W$이므로 최소성에서
$$
\|r\|^2\le \|r-tz\|^2
$$
를 얻는다. 전개하면
\begin{align*}
0&\le \|r-tz\|^2-\|r\|^2\\
&= -\overline t\,\langle r,z\rangle - t\,\langle z,r\rangle + |t|^2\|z\|^2.
\end{align*}
$z=0$이면 자명하다. $z\ne0$이면
$$
t=\frac{\langle r,z\rangle}{\|z\|^2}
$$
를 대입하여
$$
0\le -\frac{|\langle r,z\rangle|^2}{\|z\|^2}.
$$
좌변은 음이 아니고 우변은 음이 아닌 수의 음수이므로 반드시
$$
\langle r,z\rangle=0
$$
이어야 한다. $z\in W$가 임의였으므로 $r\in W^\perp$, 즉 (i)가 성립한다.

유일성은 다음과 같다. $w_1,w_2\in W$가 모두 최소거리점을 이루면 위 동치에 의해
$$
v-w_1\in W^\perp,\qquad v-w_2\in W^\perp.
$$
따라서
$$
w_1-w_2=(v-w_2)-(v-w_1)\in W^\perp
$$
이고 동시에 $w_1-w_2\in W$이다. 그러므로 $w_1-w_2\in W\cap W^\perp=\{0\}$, 즉 $w_1=w_2$.
\end{proof}

\paragraph{의미 해석.}
직교투영은 단순한 기하 그림이 아니라 ``오차를 최소로 만드는 유일한 점''입니다. 데이터 근사 문제에서 최적해가 존재하고 해석 가능한 이유가 이 정리로 완성됩니다.

\begin{example}
$\R^3$에서
$$
W=\operatorname{span}\{(1,1,0),(0,0,1)\}
$$
이라 하자. $W$의 정규직교기저를
$$
e_1=\frac1{\sqrt2}(1,1,0),\qquad e_2=(0,0,1)
$$
로 잡으면 $v=(2,1,3)$의 $W$ 위 직교투영은
\begin{align*}
P_W(v)
&=\langle v,e_1\rangle e_1+\langle v,e_2\rangle e_2\\
&=\frac3{\sqrt2}\cdot\frac1{\sqrt2}(1,1,0)+3(0,0,1)\\
&=\left(\frac32,\frac32,3\right).
\end{align*}
잔차 $v-P_W(v)=\left(\frac12,-\frac12,0\right)$는 $W$의 두 생성벡터와 모두 직교한다.
\end{example}

\subsection*{주의}
\begin{warning}
$W^\perp$는 집합론적 여집합이 아니다. 직교 조건으로 정의된 부분공간이며, 원소 개수나 포함관계와 무관하게 계산해야 한다.
\end{warning}
\begin{warning}
투영 공식을 사용할 때는 기저가 반드시 정규직교여야 한다. 일반 기저에 같은 식을 대입하면 잘못된 계수가 나온다.
\end{warning}

\subsection*{자가진단퀴즈}
\begin{enumerate}[label=\arabic*.]
\item $\R^4$에서 $W=\operatorname{span}\{(1,1,0,0),(0,1,1,0)\}$일 때 $W^\perp$의 기저를 구하라.
\item 최소거리 정리의 (ii)$\Rightarrow$(i)에서 $t=\langle r,z\rangle/\|z\|^2$를 선택하는 이유를 설명하라.
\item $P_W$가 선형사상임을 증명하고 $P_W^2=P_W$를 보이라.
\end{enumerate}

\sectiontemplate{최소제곱과 QR 분해}{관측 데이터가 과잉결정(overdetermined)일 때는 정확해 대신 오차가 가장 작은 근사해를 구합니다. 이 절에서는 직교투영 이론을 정상방정식으로 옮기고, 수치적으로 안정적인 QR 분해 방식까지 정리하겠습니다.}{\item 최소제곱 문제를 직교투영 문제로 변환할 수 있습니다.\item 정상방정식의 유도와 해의 유일성 조건을 증명할 수 있습니다.\item 열독립 행렬의 QR 분해를 구성하고 최소제곱 계산에 연결할 수 있습니다.}{\item 열공간 $\operatorname{col}(A)$\item 켤레전치 $A^*$\item 직교투영의 최소거리 성질}

\begin{definition}
$A\in M_{m,n}(K)$, $b\in K^m$라 하자. 함수
$$
f(x):=\|Ax-b\|^2\qquad (x\in K^n)
$$
를 최소로 만드는 $x$를 \emph{최소제곱해}(least squares solution)라 한다.
\end{definition}

\begin{theorem}[정상방정식]
$A\in M_{m,n}(K)$, $b\in K^m$라 하자.
\begin{enumerate}[label=(\roman*)]
\item $\hat x\in K^n$가 최소제곱해일 필요충분조건은
$$
A^*(A\hat x-b)=0
$$
이다.
\item $A$의 열벡터들이 선형독립이면 최소제곱해는 유일하고
$$
\hat x=(A^*A)^{-1}A^*b
$$
로 주어진다.
\end{enumerate}
\end{theorem}

\paragraph{증명 전략.}
핵심은 $Ax$가 항상 $\operatorname{col}(A)$에 있다는 점입니다. 따라서 최소제곱은 $b$를 열공간에 직교투영하는 문제로 바뀌며, ``잔차가 열공간과 직교''라는 조건이 곧 정상방정식으로 번역됩니다. 유일성은 $A^*A$의 양의정부호성으로 처리합니다.

\begin{proof}
$W:=\operatorname{col}(A)\subset K^m$로 두자. 임의의 $x$에 대해 $Ax\in W$이므로
$$
\min_{x\in K^n}\|Ax-b\|
=\min_{y\in W}\|y-b\|.
$$
따라서 $\hat x$가 최소제곱해라는 것은 $A\hat x$가 $b$의 $W$ 위 최소거리점이라는 뜻이다. 직교투영의 최소거리 정리에 의해 이것은
$$
b-A\hat x\in W^\perp
$$
와 동치이다. $W^\perp$ 조건은 $W$의 생성벡터인 $A$의 열벡터 $a_1,\dots,a_n$에 대해
$$
\langle b-A\hat x,a_j\rangle=0\quad(1\le j\le n)
$$
와 동치이고, 이는 행렬식으로
$$
A^*(b-A\hat x)=0
$$
즉
$$
A^*(A\hat x-b)=0
$$
와 같다. (i)가 증명되었다.

이제 열벡터 선형독립을 가정하자. $x\ne0$이면 $Ax\ne0$이므로
$$
x^*A^*Ax=\langle Ax,Ax\rangle=\|Ax\|^2>0.
$$
따라서 $A^*A$는 양의정부호이므로 가역이다. 정상방정식
$$
A^*A\hat x=A^*b
$$
는 유일해를 가지며, 그 해는
$$
\hat x=(A^*A)^{-1}A^*b
$$
이다. (i)에 의해 이 해가 곧 유일한 최소제곱해다.
\end{proof}

\paragraph{의미 해석.}
최소제곱은 미분 공식이 아니라 직교투영 구조에서 자연스럽게 나오며, 정상방정식은 그 구조를 좌표계로 옮긴 표현입니다. ``잔차가 열공간과 직교''라는 문장을 기억하면 공식 암기 없이도 전체 논리를 복구할 수 있습니다.

\begin{theorem}[QR 분해의 존재]
$A\in M_{m,n}(K)$의 열벡터 $a_1,\dots,a_n$이 선형독립이라 하자. 그러면 다음을 만족하는 행렬 $Q,R$가 존재한다.
\begin{enumerate}[label=(\roman*)]
\item $A=QR$,
\item $Q=[q_1\ \cdots\ q_n]\in M_{m,n}(K)$의 열벡터는 정규직교이다($Q^*Q=I_n$),
\item $R=(r_{ij})\in M_{n,n}(K)$는 상삼각행렬이고 $r_{jj}>0$이다.
\end{enumerate}
\end{theorem}

\paragraph{증명 전략.}
$A$의 열벡터에 Gram--Schmidt를 적용해 정규직교 열벡터 $q_j$를 만든 뒤, 각 원래 열벡터 $a_j$를 이전 $q_i$들의 선형결합으로 다시 쓰면 계수행렬이 자동으로 상삼각행렬이 됩니다.

\begin{proof}
열벡터 $a_1,\dots,a_n$에 Gram--Schmidt를 적용하여 서로 직교인 $u_1,\dots,u_n$을 얻고
$$
q_j:=\frac{u_j}{\|u_j\|}\quad(1\le j\le n)
$$
로 둔다. 선형독립성 때문에 각 $u_j\ne0$이므로 정의가 가능하며, $q_1,\dots,q_n$은 정규직교이다.

Gram--Schmidt 구성에서 각 $j$에 대해
$$
a_j=\sum_{i=1}^j r_{ij}q_i
$$
꼴로 쓸 수 있고,
$$
r_{ij}=\langle a_j,q_i\rangle\ (i<j),\qquad r_{jj}=\|u_j\|>0,\qquad r_{ij}=0\ (i>j)
$$
로 둘 수 있다. 이제
$$
Q:=[q_1\ \cdots\ q_n],\qquad R:=(r_{ij})
$$
라 두면 $j$번째 열식이 정확히 $a_j=\sum_{i=1}^j r_{ij}q_i$이므로 행렬식으로
$$
A=QR
$$
가 된다. $q_i$들이 정규직교이므로 $Q^*Q=I_n$이며, $R$은 정의상 상삼각이고 대각원소가 양수다.
\end{proof}

\paragraph{의미 해석.}
QR 분해는 ``직교화 + 좌표변환''을 한 번에 묶은 형태입니다. 최소제곱 계산에서 $A^*A$를 직접 만들지 않아도 되어 수치적으로 더 안정적인 절차를 제공합니다.

\begin{example}
세 점 $(0,1),(1,2),(2,2)$를 직선 $y=\beta_0+\beta_1 t$로 근사하자. 
$$
A=
\begin{bmatrix}
1&0\\
1&1\\
1&2
\end{bmatrix},\qquad
b=
\begin{bmatrix}
1\\2\\2
\end{bmatrix}.
$$
정상방정식은
$$
A^TA=
\begin{bmatrix}
3&3\\
3&5
\end{bmatrix},\qquad
A^Tb=
\begin{bmatrix}
5\\6
\end{bmatrix}
$$
이므로
$$
\hat x=
\begin{bmatrix}
\beta_0\\\beta_1
\end{bmatrix}
=
\begin{bmatrix}
7/6\\1/2
\end{bmatrix}.
$$
따라서 최소제곱 직선은
$$
y=\frac76+\frac12 t
$$
이다. 잔차 벡터 $r=b-A\hat x=(-1/6,\,1/3,\,-1/6)^T$는 $A$의 두 열벡터와 직교한다.
\end{example}

\subsection*{주의}
\begin{warning}
정상방정식은 이론적으로 간단하지만 수치계산에서는 조건수를 제곱시키는 경향이 있다. 실제 계산에서는 QR 분해 방식이 더 안정적일 수 있다.
\end{warning}
\begin{warning}
열독립이 깨지면 최소제곱해는 존재하더라도 유일하지 않을 수 있다. ``해의 존재''와 ``해의 유일성''을 반드시 분리해서 판단해야 한다.
\end{warning}

\subsection*{자가진단퀴즈}
\begin{enumerate}[label=\arabic*.]
\item 최소제곱 문제를 ``$b$를 $\operatorname{col}(A)$에 투영하는 문제''로 바꾸는 과정을 식으로 정리하라.
\item $A=\begin{bmatrix}1&1\\1&-1\\1&0\end{bmatrix}$에 대해 QR 분해를 구성하고 $Q^*Q=I$를 확인하라.
\item 열독립인 경우 $A^*A$가 양의정부호임을 직접 증명하라.
\end{enumerate}

\chapterapplicationtemplate{센서 보정 데이터를 선형모형 $b\approx Ax$로 설정하면 관측 잡음 때문에 정확한 해가 없을 때가 많습니다. 이 장의 결과를 사용하면, 먼저 열공간에 대한 직교투영으로 오차가 최소인 예측값 $A\hat x$를 정의하고, 이어서 QR 분해로 안정적으로 $\hat x$를 계산할 수 있습니다. 특히 잔차 $r=b-A\hat x$가 모델 공간과 직교한다는 사실은 ``현재 모형이 더 이상 설명하지 못하는 정보''를 정량적으로 분리하는 기준이 됩니다.}

\chaptersummarytemplate

\section*{종합 연습문제}

\subsection*{연습문제 A (기초 확인)}
\begin{enumerate}[label=\textbf{A\arabic*.}]
\item $\C^2$에서 $\langle x,y\rangle=2x_1\overline{y_1}+x_2\overline{y_2}$가 내적인지 판정하라.
\item $x=(1,-1,2)$, $y=(2,1,0)$에 대해 $\langle x,y\rangle$, $\|x\|$, $\|y\|$를 계산하고 Cauchy--Schwarz를 수치로 확인하라.
\item 벡터열 $(1,1,0),(1,0,1),(0,1,1)$에 Gram--Schmidt를 적용하여 정규직교집합을 구하라.
\item $W=\operatorname{span}\{(1,1,0)\}\subset\R^3$일 때 $v=(2,0,1)$의 $P_W(v)$를 구하라.
\item $W=\operatorname{span}\{(1,0,1),(0,1,1)\}\subset\R^3$의 직교여공간 $W^\perp$의 기저를 구하라.
\item 
$$
A=
\begin{bmatrix}
1&0\\
1&1\\
1&2
\end{bmatrix},\quad
b=
\begin{bmatrix}
1\\0\\2
\end{bmatrix}
$$
에 대해 정상방정식을 세우고 최소제곱해를 구하라.
\end{enumerate}

\subsection*{연습문제 B (개념 연결)}
\begin{enumerate}[label=\textbf{B\arabic*.}]
\item $W$의 정규직교기저가 주어졌을 때 $P_W$가 선형사상임을 증명하고 $\operatorname{Im}(P_W)=W$, $\ker(P_W)=W^\perp$를 보여라.
\item 직교분해 $v=w+u$ ($w\in W$, $u\in W^\perp$)가 유일함을 교집합 조건으로 다시 증명하라.
\item 열독립 행렬 $A$에 대해 $A=QR$이면 최소제곱 문제가 $R\hat x=Q^*b$의 해법으로 바뀜을 유도하라.
\item 최소제곱 잔차 $r=b-A\hat x$가 0이면 어떤 의미인지, 그리고 이 조건이 언제 가능한지 열공간 관점에서 설명하라.
\end{enumerate}

\subsection*{연습문제 C (증명/종합)}
\begin{enumerate}[label=\textbf{C\arabic*.}]
\item 유한차원 내적공간의 정규직교기저 $\{e_1,\dots,e_n\}$에 대해, 임의의 $x,y\in V$가
$$
\langle x,y\rangle=\sum_{i=1}^n \langle x,e_i\rangle\overline{\langle y,e_i\rangle}
$$
를 만족함을 증명하라.
\item $T:V\to V$가 $\langle Tx,Ty\rangle=\langle x,y\rangle$를 모든 $x,y$에 대해 만족한다고 하자. $V$가 유한차원일 때 $T$가 가역임을 증명하고, $T^{-1}$을 내적 조건으로 표현하라.
\end{enumerate}